\section{Proof Sketches and Boundary Cases}

We provide proof sketches for the main results and identify boundary cases that
delineate the scope of the phase-transition claim. All arguments are internal to
the OC/ETS formalism and rely exclusively on admissibility, thresholds, and cycle
structure.

\subsection{Proof Sketch: Identity Loss Transition}

Identity loss occurs when either
(i) the identity-supporting subspace $\Lambda_I$ intersects the boundary
$\partial\Omega_{\mathrm{EA}}$, or
(ii) its supporting cycle set $C_I$ becomes empty.

In case (i), admissibility is lost by definition: states arbitrarily close to
$\partial\Omega_{\mathrm{EA}}$ violate threshold constraints, so $\Lambda_I$ can
no longer function as an admissible invariant subspace. In case (ii), invariance
is lost: without cycles, there exist no closed trajectories capable of sustaining
identity-equivalent motion. In both cases, the set of admissible invariant
trajectories undergoes a discontinuous change while the admissible state space
$\Omega_{\mathrm{EA}}$ itself may remain non-empty. This satisfies the OC criteria
for a phase transition.

\subsection{Proof Sketch: Irreversibility}

Irreversibility follows from closure of the admissible dynamics. If cycles in the
identity-relevant class cannot be regenerated without introducing new admissible
structure, then once $C_I = \varnothing$ there exists no internal mechanism by
which an identity-supporting subspace can re-emerge. Any putative recovery would
require extension of $A_{\mathrm{EA}}$, $P_{\mathrm{EA}}$, or
$\Theta_{\mathrm{EA}}$, which lies outside the closed system under
consideration.

\subsection{Boundary Case: Inertia without Identity}

A critical boundary case arises when $\Omega_{\mathrm{EA}} \neq \varnothing$ and
$k_{\mathrm{EA}}$ remains positive, yet $C_I = \varnothing$. In this regime the
enterprise exhibits admissible operation and may persist in prolonged inertia,
but no identity in the formal sense defined above. This demonstrates that
identity loss is neither equivalent to collapse nor necessarily accompanied by
immediate organizational failure.

\subsection{Boundary Case: Apparent Identity Persistence}

Another boundary case occurs when trajectories remain close to the former
$\Lambda_I$ while approaching $\partial\Omega_{\mathrm{EA}}$. Such behavior may
appear as identity persistence under coarse observation; however, formally
$\Lambda_I$ is no longer bounded away from the boundary and therefore fails to be
admissible. This clarifies that identity loss can precede observable collapse and
may remain masked until threshold crossing is completed.

These boundary cases demonstrate that identity loss defines a distinct structural
phase, separable both from full collapse and from ordinary operational
fluctuation within $\Omega_{\mathrm{EA}}$.
